% ============================================================
% UNIDAD 1: DEL PANEL ANALÓGICO A LA CABINA DE CRISTAL
% ============================================================
\clearpage
\section*{UNIDAD 1: LA CABINA DE CRISTAL}
\subsection*{Del panel analógico a los instrumentos integrados}

\begin{recuadro}
\textbf{Objetivo de la unidad:} conocer la clasificación y agrupación de los instrumentos de
vuelo, la evolución hacia los sistemas electrónicos de vuelo (EFIS), las arquitecturas de los
instrumentos integrados y el marco de normas y certificación en el que se instalan.
\end{recuadro}

% ------------------------------------------------------------
\section*{1. LOS INSTRUMENTOS CLÁSICOS Y LA T BÁSICA}

Los instrumentos de una aeronave se clasifican en \textbf{instrumentos de vuelo} (actitud,
velocidad, altitud, rumbo), \textbf{de navegación} (VOR, ILS, GPS, ADF), \textbf{de motor}
(revoluciones, presiones, temperaturas, combustible) y \textbf{de sistemas} (eléctrico,
hidráulico, presurización, tren, flaps).

Durante décadas los instrumentos de vuelo fueron mecánicos o electromecánicos, agrupados en la
\textbf{T básica}: actitud en el centro, velocidad a la izquierda, altitud a la derecha y rumbo
debajo, para que la mirada del piloto recorra siempre el mismo camino. Con el indicador de viraje
y el variómetro se forma el clásico panel de ``seis instrumentos''.

\begin{center}
\begin{tikzpicture}[every node/.style={font=\footnotesize, align=center}]
  \draw[ppAzul!25, line width=14pt] (0,2.6) -- (6,2.6);
  \draw[ppAzul!25, line width=14pt] (3,2.6) -- (3,0);
  \draw[thick, fill=gray!10] (0,2.6) circle (1.15);
  \node[font=\footnotesize\bfseries] at (0,2.8) {Velocímetro};
  \node[gray] at (0,2.0) {\tiny pitot-estático};
  \draw[thick, fill=gray!10] (3,2.6) circle (1.15);
  \node[font=\footnotesize\bfseries] at (3,2.8) {Horizonte\\ artificial};
  \node[gray, font=\tiny] at (3,2.0) {giróscopo\\ (vacío)};
  \draw[thick, fill=gray!10] (6,2.6) circle (1.15);
  \node at (6,2.8) {\textbf{Altímetro}};
  \node[gray] at (6,2.0) {\tiny estático};
  \draw[thick, fill=gray!10] (0,0) circle (1.15);
  \node[font=\footnotesize\bfseries] at (0,0.2) {Coordinador\\ de viraje};
  \node[gray, font=\tiny] at (0,-0.6) {giróscopo\\ (eléctrico)};
  \draw[thick, fill=gray!10] (3,0) circle (1.15);
  \node[font=\footnotesize\bfseries] at (3,0.2) {Indicador\\ de rumbo};
  \node[gray, font=\tiny] at (3,-0.6) {giróscopo\\ (vacío)};
  \draw[thick, fill=gray!10] (6,0) circle (1.15);
  \node at (6,0.2) {\textbf{Variómetro}};
  \node[gray] at (6,-0.6) {\tiny estático};
  \node[ppAzul, right] at (7.4,1.3) {\textbf{T básica}};
\end{tikzpicture}
\end{center}

\paragraph*{Problemas del panel analógico.} Instrumentos de vacío con giróscopos que se desgastan;
bombas de vacío que fallan sin aviso; mucho peso y cableado; poca información integrada; cada
instrumento con su propio mantenimiento y calibración.

% ------------------------------------------------------------
\section*{2. LOS SISTEMAS ELECTRÓNICOS DE VUELO (EFIS)}

En los años ochenta los aviones de transporte (Boeing 757/767, Airbus A310) reemplazaron los
instrumentos de actitud y de rumbo por pantallas de tubos de rayos catódicos: el \textbf{EFIS}
(\textit{Electronic Flight Instrument System}). Luego llegaron las pantallas de cristal líquido y la
integración total. En la aviación general, el \textbf{Garmin G1000} (2004) llevó la
\textbf{cabina de cristal} (\textit{glass cockpit}) a los aviones livianos.

\begin{center}
\small
\begin{tabular}{|p{0.18\linewidth}|p{0.70\linewidth}|}
\hline
\textbf{Pantalla} & \textbf{Función} \\ \hline
\textbf{PFD} & \textit{Primary Flight Display}: reúne la T básica en una sola pantalla (actitud,
velocidad, altitud, velocidad vertical, rumbo) más la navegación, el piloto automático y las
alertas. \\ \hline
\textbf{MFD} & \textit{Multi-Function Display}: mapa móvil, motor, tránsito, meteorología, terreno,
cartas, listas de chequeo. \\ \hline
\textbf{ND} & \textit{Navigation Display}: en aviones de transporte, la pantalla de navegación. \\
\hline
\textbf{EICAS / ECAM} & Pantallas de motor y de alertas de la tripulación (Boeing y Airbus). \\
\hline
\textbf{EIS} & \textit{Engine Indication System}: indicación de motor en aviación general. \\
\hline
\end{tabular}
\end{center}

\paragraph*{Ventajas.} Más información, mejor integrada y presentada de forma más intuitiva;
menos peso; mayor confiabilidad (sin giróscopos mecánicos ni bombas de vacío); autodiagnóstico;
actualización por software; menos cableado gracias a los buses de datos.

\paragraph*{Desafíos.} Dependen totalmente del sistema eléctrico; exigen capacitación de pilotos y
técnicos; las bases de datos y el software deben mantenerse actualizados; una falla puede afectar
varias funciones a la vez. Un estudio de la NTSB de 2010 sobre aviones livianos con cabina de
cristal concluyó que la tecnología por sí sola no mejoraba la seguridad si no venía acompañada de
\textbf{capacitación específica}.

% ------------------------------------------------------------
\section*{3. ARQUITECTURAS}

\begin{itemize}
  \item \textbf{Instrumentos autónomos (\textit{standalone}):} una pantalla con sus sensores
        internos o propios reemplaza a uno o dos instrumentos (G5, GI 275, Aspen).
  \item \textbf{Sistema de pantallas con sensores remotos:} las pantallas reciben los datos de un
        ADAHRS (computadora de datos de aire y referencia de actitud) y otros sensores, y se
        comunican con la radio, el GPS y el piloto automático existentes (G3X Touch, G500/G600
        TXi).
  \item \textbf{Sistema totalmente integrado:} pantallas, sensores, radios, transpondedor, audio y
        piloto automático forman un solo sistema conectado por buses de alta velocidad, con
        unidades de procesamiento centrales (G1000, G3000, G5000, y en los aviones de transporte la
        aviónica modular integrada, IMA).
  \item \textbf{Pantalla con radios integradas:} la tendencia más reciente en aviación general es
        que la pantalla incorpore también el navegador GPS, las radios de comunicaciones y
        navegación y el panel de audio, eliminando la columna central de equipos (Garmin AXIS).
\end{itemize}

\paragraph*{Instrumentos de reserva.} Toda cabina de cristal mantiene \textbf{instrumentos de
reserva} (\textit{standby}) independientes, con su propia batería, para volar si fallan las
pantallas principales: antes eran un horizonte, un altímetro y un velocímetro mecánicos; hoy suele
ser un instrumento electrónico (por ejemplo, un G5 o un GI 275).

% ------------------------------------------------------------
\section*{4. NORMAS Y CERTIFICACIÓN}

\begin{center}
\small
\begin{tabular}{|p{0.26\linewidth}|p{0.62\linewidth}|}
\hline
\textbf{Norma o documento} & \textbf{Qué regula} \\ \hline
TSO (\textit{Technical Standard Order}) & Estándar mínimo de desempeño de un equipo (por ejemplo,
TSO-C113 para pantallas multipropósito). \\ \hline
RTCA DO-160 & Ensayos ambientales: temperatura, vibración, humedad, interferencia, rayos. \\
\hline
RTCA DO-178C / DO-254 & Desarrollo del software y del hardware electrónico complejo. \\ \hline
FAA AC 23.1311-1 & Instalación de pantallas electrónicas en aviones de categoría normal:
ubicación, campo visual, colores, alertas. \\ \hline
FAA AC 43.13-1B y 43.13-2 & Prácticas aceptables de mantenimiento, cableado e instalación de
equipos. \\ \hline
STC (certificado tipo suplementario) & Aprobación de una modificación de diseño. Un
\textbf{AML-STC} cubre una lista de modelos de aeronave. \\ \hline
\end{tabular}
\end{center}

\paragraph*{Aeronaves certificadas, experimentales y LSA.} En una aeronave \textbf{certificada}
solo se pueden instalar equipos con aprobación (TSO) y mediante datos aprobados (STC o aprobación
de campo), y la instalación se documenta como alteración ante la autoridad. En las aeronaves
\textbf{experimentales} (de construcción amateur) y \textbf{LSA} hay más libertad y se usan
versiones experimentales de los equipos, más económicas. Por eso muchos fabricantes ofrecen dos
versiones del mismo equipo. En la Argentina, la instalación en una aeronave certificada requiere
que el STC extranjero esté \textbf{convalidado o aceptado por la ANAC} y que la realice un taller
habilitado.

% ------------------------------------------------------------
\section*{5. ACTIVIDADES DE LA UNIDAD}
\begin{enumerate}
  \item Dibujar la T básica e indicar qué sistema alimenta a cada instrumento (pitot, estática,
        vacío, eléctrico).
  \item Comparar una foto de un panel analógico y una de un panel con G1000: ubicar en el PFD cada
        uno de los seis instrumentos clásicos.
  \item Investigar qué aeronaves de la escuela o del aeródromo más cercano tienen cabina de
        cristal y con qué equipos.
  \item Explicar la diferencia entre un instrumento autónomo, un sistema con sensores remotos y un
        sistema integrado, con un ejemplo de cada uno.
  \item Buscar en la lista de modelos (AML) de un STC de Garmin si incluye a un Cessna 150 o un
        Piper PA-28.
\end{enumerate}
